\paragraph{Bernoulli--\texorpdfstring{$\zeta(2r)$}{zeta(2r)} 塔的稀疏刚性、primitive 长度 \texorpdfstring{$2$}{2} 原子的同余局域化与 escort 的一比特冻结}
在上段关于 Wedderburn 块谱、fold-gauge 常数层级、真实输入核 primitive 原子以及 bin-fold 两态一致渐近的结论之外，还可继续抽取四条更细但仍完全可审计的后果：其一是 \(\log C_m\) 的任意固定算术子列已经足以决定整条 Bernoulli--\(\zeta(2r)\) 塔；其二是真实输入核中可剥离的 primitive 长度 \(2\) 原子在任何 root-of-unity 长度滤波下都只落在唯一的同余类；其三是 bin-fold 的 escort 热力学在主阶上完全退化为末位比特的显式 Gibbs 偏置；其四是相应的 min-entropy 率从 \(a=0\) 起即已冻结到 \(\log\varphi\)。

\begin{theorem}[Bernoulli--\texorpdfstring{$\zeta(2r)$}{zeta(2r)} 塔的算术子列刚性]\label{thm:gut-logCm-arithmetic-subsequence-rigidity}
记
\[
q:=\frac{\varphi}{2}\in(0,1),
\qquad
A_r:=\frac{B_{2r}}{r(2r-1)}\bigl(\varphi^{-1}+\varphi^{2r-3}\bigr)
\qquad (r\ge 1),
\]
其中 \(B_{2r}\) 为定理 \ref{thm:fold-bin-gauge-constant-stirling-bernoulli-hierarchy} 中出现的 Bernoulli 系数。则对任意固定算术子列
\[
m=c+dn\qquad(d\ge 1,\ 0\le c<d),
\]
仅凭充分大 \(n\) 上的稀疏样本
\[
\bigl\{\log C_{c+dn}\bigr\}_{n\gg 1}
\]
就可唯一恢复全部 \(\{A_r\}_{r\ge 1}\)，从而唯一恢复全部 \(\{B_{2r}\}_{r\ge 1}\) 与 \(\{\zeta(2r)\}_{r\ge 1}\)。

等价地，若另一列 \(\widetilde C_m\) 满足同型渐近展开，且
\[
\log C_{c+dn}=\log \widetilde C_{c+dn}\qquad(n\gg 1),
\]
则两列对应的全部 Bernoulli 系数与全部偶值 \(\zeta(2r)\) 必逐阶一致。
\leanverified{paper\_gut\_logCm\_arithmetic\_subsequence\_rigidity\_seeds}
\leanverified{paper\_gut\_logCm\_arithmetic\_subsequence\_rigidity\_package}
\leanverified{paper\_gut\_logCm\_arithmetic\_subsequence\_rigidity}
\end{theorem}
\begin{proof}
由定理 \ref{thm:fold-bin-gauge-constant-stirling-bernoulli-hierarchy}，对任意固定 \(R\ge 1\) 有
\[
\log C_m
=\log(2\pi)+\sum_{r=1}^{R}A_r q^{(2r-1)m}+O\!\bigl(q^{(2R+1)m}\bigr).
\]
沿算术子列 \(m=c+dn\) 代入，得到
\[
\log C_{c+dn}
=\log(2\pi)+\sum_{r=1}^{R}A_r q^{(2r-1)c}(q^d)^{(2r-1)n}
+O\!\bigl((q^d)^{(2R+1)n}\bigr).
\]
因此在子列参数 \(n\) 上，仍然出现严格分离的指数尺度
\(
(q^d)^{(2r-1)n}
\)
（\(r\ge 1\)）。

设另一列 \(\widetilde C_m\) 在同一算术子列上最终与 \(C_m\) 相同，并令 \(\widetilde A_r\) 为其对应系数。若存在最小 \(r_0\) 使 \(\Delta_{r_0}:=A_{r_0}-\widetilde A_{r_0}\neq 0\)，则两列之差满足
\[
0
=
\Delta_{r_0}q^{(2r_0-1)c}(q^d)^{(2r_0-1)n}
+O\!\bigl((q^d)^{(2r_0+1)n}\bigr).
\]
两边同除以 \(q^{(2r_0-1)c}(q^d)^{(2r_0-1)n}\)，并令 \(n\to\infty\)，由于 \(0<q^d<1\)，余项趋于 \(0\)，从而得到 \(\Delta_{r_0}=0\)，矛盾。故所有系数逐阶一致。

于是 \(\{A_r\}_{r\ge 1}\) 由任意固定算术子列唯一决定。再由
\[
A_r=\frac{B_{2r}}{r(2r-1)}\bigl(\varphi^{-1}+\varphi^{2r-3}\bigr)
\]
可唯一恢复全部 \(B_{2r}\)，并由标准恒等式
\[
B_{2r}=(-1)^{r-1}\frac{2(2r)!}{(2\pi)^{2r}}\zeta(2r)
\]
唯一恢复全部 \(\zeta(2r)\)。证毕。
\end{proof}

\begin{theorem}[真实输入核 primitive 长度 \texorpdfstring{$2$}{2} 原子的同余局域化]\label{thm:gut-real-input-primitive-2-atom-congruence-localization}
\leanverified{paper\_gut\_real\_input\_primitive\_2\_atom\_congruence\_localization}
沿用附录 \ref{app:real-input-40-zeta-u} 的记号。真实输入 \(40\) 状态核的一参数加权 family \(M(u)\) 满足
\[
\Delta(z,u)=\det(I-zM(u))=(1-uz^2)\,F(z,u),
\]
并且 primitive 生成函数存在精确原子剥离
\[
P(z,u)=uz^2+P_{\mathrm{core}}(z,u)
\]
（命题 \ref{prop:atomic-2-orbit-split-logM}）。

对任意模数 \(q\ge 2\) 与剩余类 \(a\in\ZZ/q\ZZ\)，定义 primitive 长度同余滤波
\[
P^{[a,q]}(z,u):=\frac1q\sum_{j=0}^{q-1}\zeta_q^{-aj}P(\zeta_q^j z,u),
\qquad \zeta_q:=e^{2\pi i/q},
\]
以及
\[
P_{\mathrm{core}}^{[a,q]}(z,u):=\frac1q\sum_{j=0}^{q-1}\zeta_q^{-aj}P_{\mathrm{core}}(\zeta_q^j z,u).
\]
则有完全显式的局域化公式
\[
\boxed{\;
P^{[a,q]}(z,u)
=
\mathbf 1_{\{a\equiv 2\ (\mathrm{mod}\ q)\}}\,u z^2
+P_{\mathrm{core}}^{[a,q]}(z,u).\;}
\]
换言之，该 primitive 长度 \(2\) 原子在所有长度同余类分解中，只会落在唯一的剩余类 \(2\bmod q\)；特别地，当 \(q=2\) 时它全部落在偶长通道，对奇长通道贡献严格为零。
\end{theorem}
\begin{proof}
把命题 \ref{prop:atomic-2-orbit-split-logM} 的分解
\(
P(z,u)=uz^2+P_{\mathrm{core}}(z,u)
\)
代入滤波定义，得
\[
P^{[a,q]}(z,u)
=\frac1q\sum_{j=0}^{q-1}\zeta_q^{-aj}\bigl(u\zeta_q^{2j}z^2+P_{\mathrm{core}}(\zeta_q^j z,u)\bigr).
\]
第二项正是 \(P_{\mathrm{core}}^{[a,q]}(z,u)\)。第一项化为
\[
u z^2\cdot \frac1q\sum_{j=0}^{q-1}\zeta_q^{(2-a)j}
=u z^2\,\mathbf 1_{\{a\equiv 2\ (\mathrm{mod}\ q)\}},
\]
其中最后一步使用根单位正交关系。故结论成立。证毕。
\end{proof}

\begin{theorem}[bin-fold escort 族的一比特 Gibbs 化与从 \texorpdfstring{$a=0$}{a=0} 即开始的冻结]\label{thm:gut-foldbin-escort-one-bit-gibbs-freezing}
沿用 bin-fold 的纤维记号 \(d_m^{\mathrm{bin}}(w)\) 与稳定词层 \(X_m\)。对任意固定 \(a\ge 0\)，定义 escort 概率
\[
\pi_{m,a}^{\mathrm{bin}}(w):=
\frac{\bigl(d_m^{\mathrm{bin}}(w)\bigr)^a}{\sum_{v\in X_m}\bigl(d_m^{\mathrm{bin}}(v)\bigr)^a},
\qquad
w\in X_m,
\]
并令
\[
W_m^{(a)}=(W^{(a)}_{m,1},\dots,W^{(a)}_{m,m})\sim \pi_{m,a}^{\mathrm{bin}}.
\]
则有精确极限律
\[
\mathbb{P}\bigl(W^{(a)}_{m,m}=1\bigr)\longrightarrow \frac{1}{1+\varphi^{a+1}},
\qquad
\mathbb{P}\bigl(W^{(a)}_{m,m}=0\bigr)\longrightarrow \frac{\varphi^{a+1}}{1+\varphi^{a+1}}.
\]
并且
\[
\log d_m^{\mathrm{bin}}\!\bigl(W_m^{(a)}\bigr)-m\log\!\Bigl(\frac{2}{\varphi}\Bigr)
\]
在分布上收敛到两点分布
\[
\begin{cases}
0,& \text{以概率 }\displaystyle \frac{\varphi^{a+1}}{1+\varphi^{a+1}},\\[8pt]
-\log\varphi,& \text{以概率 }\displaystyle \frac{1}{1+\varphi^{a+1}}.
\end{cases}
\]
此外，对 escort 的 min-entropy 有
\[
-\frac1m\log\!\Bigl(\max_{w\in X_m}\pi_{m,a}^{\mathrm{bin}}(w)\Bigr)\longrightarrow \log\varphi
\qquad (a\ge 0).
\]
因此，bin-fold 的 escort 热力学在主阶上完全退化为末位比特的一比特 Gibbs 偏置，并且其 min-entropy 率从 \(a=0\) 起即已冻结到 \(\log\varphi\)。
\leanverified{paper\_gut\_foldbin\_escort\_one\_bit\_gibbs\_freezing\_seeds}
\leanverified{paper\_gut\_foldbin\_escort\_one\_bit\_gibbs\_freezing}
\end{theorem}
\begin{proof}
末位极限律正是命题 \ref{prop:fold-bin-escort-lastbit} 在参数 \(q=a\) 时的结论，即
\[
\mathbb{P}\bigl(W^{(a)}_{m,m}=1\bigr)\to \theta(a):=\frac{1}{1+\varphi^{a+1}},
\qquad
\mathbb{P}\bigl(W^{(a)}_{m,m}=0\bigr)\to 1-\theta(a).
\]

由定理 \ref{thm:fold-bin-two-state-asymptotic} 的一致展开，
\[
\log d_m^{\mathrm{bin}}(w)
=m\log\!\Bigl(\frac{2}{\varphi}\Bigr)-w_m\log\varphi+O\!\Bigl(\Bigl(\frac{\varphi}{2}\Bigr)^m\Bigr)
\qquad (w\in X_m).
\]
故
\[
\log d_m^{\mathrm{bin}}\!\bigl(W_m^{(a)}\bigr)-m\log\!\Bigl(\frac{2}{\varphi}\Bigr)
=-W^{(a)}_{m,m}\log\varphi+o_{\mathbb P}(1),
\]
与上式的末位分布合并，即得所示两点极限分布；这也可视为推论 \ref{cor:fold-bin-escort-two-scale-residual} 经对数映射后的直接后果。

最后证明 min-entropy 率。若 \(a=0\)，则 \(\pi_{m,0}^{\mathrm{bin}}\) 就是 \(X_m\) 上的均匀分布，故
\[
\max_{w\in X_m}\pi_{m,0}^{\mathrm{bin}}(w)=\frac{1}{|X_m|}=\frac{1}{F_{m+2}},
\]
从而
\[
-\frac1m\log\max_{w}\pi_{m,0}^{\mathrm{bin}}(w)
=\frac{1}{m}\log F_{m+2}\longrightarrow \log\varphi.
\]
若 \(a>0\)，则由定理 \ref{thm:fold-bin-two-state-asymptotic} 与命题 \ref{prop:fold-bin-power-sum-asymptotic}，
\[
\pi_{m,a}^{\mathrm{bin}}(w)
=\frac{\varphi^{-a w_m}}{F_{m+1}+\varphi^{-a}F_m}\Bigl(1+O\!\Bigl(\Bigl(\frac{\varphi}{2}\Bigr)^m\Bigr)\Bigr)
\]
在 \(w\in X_m\) 上一致成立。于是 \(w_m=0\) 扇区的点质量在主阶上最大，且
\[
\max_{w\in X_m}\pi_{m,a}^{\mathrm{bin}}(w)
=\frac{1}{F_{m+1}+\varphi^{-a}F_m}\bigl(1+o(1)\bigr).
\]
再用 \(F_{m+1}+\varphi^{-a}F_m=\Theta(\varphi^m)\) 即得
\[
-\frac1m\log\max_{w\in X_m}\pi_{m,a}^{\mathrm{bin}}(w)\to \log\varphi.
\]
证毕。
\end{proof}

\begin{corollary}[bin-fold escort 族的精确热力学一阶矩与二阶矩]\label{cor:gut-foldbin-escort-logfiber-first-second-moments}
沿用定理 \ref{thm:gut-foldbin-escort-one-bit-gibbs-freezing} 的记号。则有
\[
\mathbb E_{\pi_{m,a}^{\mathrm{bin}}}\!\left[\log d_m^{\mathrm{bin}}\!\bigl(W_m^{(a)}\bigr)\right]
=
m\log\!\Bigl(\frac{2}{\varphi}\Bigr)
-\frac{\log\varphi}{1+\varphi^{a+1}}
+O\!\Bigl(\Bigl(\frac{\varphi}{2}\Bigr)^m\Bigr),
\]
并且
\[
\operatorname{Var}_{\pi_{m,a}^{\mathrm{bin}}}\!\left(\log d_m^{\mathrm{bin}}\!\bigl(W_m^{(a)}\bigr)\right)
=
\frac{\varphi^{a+1}}{(1+\varphi^{a+1})^2}(\log\varphi)^2
+o(1).
\]
\leanverified{paper\_gut\_foldbin\_escort\_logfiber\_first\_second\_moments}
\end{corollary}
\begin{proof}
一阶矩展开正是命题 \ref{prop:fold-bin-escort-lastbit} 在 \(q=a\) 时对
\[
\kappa_m^{\mathrm{bin}}(a)
=
\mathbb E_{\pi_{m,a}^{\mathrm{bin}}}\!\left[\log d_m^{\mathrm{bin}}\!\bigl(W_m^{(a)}\bigr)\right]
\]
给出的显式公式。

对二阶矩，由定理 \ref{thm:gut-foldbin-escort-one-bit-gibbs-freezing}，
\[
\log d_m^{\mathrm{bin}}\!\bigl(W_m^{(a)}\bigr)-m\log\!\Bigl(\frac{2}{\varphi}\Bigr)
\]
收敛到取值 \(0\) 与 \(-\log\varphi\) 的两点分布，对应概率分别为
\[
\frac{\varphi^{a+1}}{1+\varphi^{a+1}},
\qquad
\frac{1}{1+\varphi^{a+1}}.
\]
该两点分布的方差恰为
\[
\frac{\varphi^{a+1}}{(1+\varphi^{a+1})^2}(\log\varphi)^2,
\]
故所示极限成立。证毕。
\end{proof}

这四条结果把此前分散的四个接口压缩到同一条可审计链上：任意固定算术子列 \(\{\log C_{c+dn}\}\) 已足以恢复整条 Bernoulli--\(\zeta(2r)\) 塔；真实输入核中可剥离的 primitive 长度 \(2\) 原子在任意长度同余类分解中都只支撑于唯一剩余类；bin-fold escort 族在主阶上完全退化为末位比特的显式 Gibbs 偏置；而相应的 min-entropy 率不需经过任何正温度阈值，已从 \(a=0\) 起固定为 \(\log\varphi\)。
